\documentclass{rapportECL}
\usepackage[utf8]{inputenc}
\usepackage[T1]{fontenc}
\usepackage[french]{babel}

% Gestion des marges pour A4 vertical
\usepackage[top=2.5cm, bottom=2.5cm, left=2.5cm, right=2.5cm]{geometry}

\usepackage{graphicx}
\usepackage{xcolor}
\usepackage{tikz}
\usetikzlibrary{arrows.meta, positioning, calc}

% Définition des styles pour les blocs, sommes, etc.
\tikzset{
    block/.style = {draw, rectangle, minimum height=2.5em, minimum width=4em, align=center, thick},
    sum/.style = {draw, circle, inner sep=2pt, thick},
    input/.style = {coordinate},
    output/.style = {coordinate},
    line/.style = {draw, -Latex, thick}
}

% Informations du rapport
\titre{Rendu Jumeaux numériques}\\ 
% \large UE : Dimensionnement des structures (ENISE-5GM)}
\eleves{Kevin \textsc{TONGUE}}
\enseignant{Enseignant : Renaud \textsc{FERRIER}}

\begin{document}
%----------- Initialisation -------------------
        
\fairemarges %Afficher les marges
\fairepagedegarde %Créer la page de garde
% \tabledematieres %Créer la table de matières

\vspace{1cm}

\section*{4.1 Travail préparatoire}
\subsection*{4.1.1 Étude d'un système et de son asservissement}

Le système est composé d'un moteur piloté en courant $i$, avec une fonction de transfert $F_m$ donnant le couple $C_m = F_m i$. La transmission a une fonction de transfert $F_t$ telle que la force $f_A = F_t C_m$. La poutre est représentée par une fonction de transfert $F$ ayant en entrée $f_A$ et en sortie le vecteur de déplacements $u$. Des projecteurs $\Pi_A$ et $\Pi_B$ extraient $u_A$ et $u_B$ de $u$.

\subsubsection*{Schéma blocs en boucle ouverte}

\begin{figure}[!htbp]
    \centering
    \begin{tikzpicture}[auto, node distance=2cm]
        % Noeuds de la chaîne principale
        \node [input] (input) {};
        \node [block, right=1cm of input] (F0) {$F_0$};
        \node [block, right=1cm of F0] (Ft) {$F_t$};
        \node [block, right=1cm of Ft] (F) {$F$};
        \node [coordinate, right=1cm of F] (split) {};
        
        % Sorties parallèles
        \node [block, above right=0.2cm and 1.5cm of split] (PiB) {$\Pi_B$};
        \node [block, below right=0.2cm and 1.5cm of split] (PiA) {$\Pi_A$};
        \node [output, right=1cm of PiB] (outB) {};
        \node [output, right=1cm of PiA] (outA) {};

        % Liaisons
        \draw [line] (input) -- node {$I$} (F0);
        \draw [line] (F0) -- node {$C_m$} (Ft);
        \draw [line] (Ft) -- node {$f_A$} (F);
        \draw [line] (F) -- node {$U$} (split);
        
        % Embranchement vers les projecteurs
        \draw [line] (split) |- (PiB);
        \draw [line] (split) |- (PiA);
        \draw [line] (PiB) -- node {$U_B$} (outB);
        \draw [line] (PiA) -- node {$U_A$} (outA);
    \end{tikzpicture}
    \caption{Schéma bloc du système en boucle ouverte}
    \label{fig:sys_boucle_ouverte}
\end{figure}

Le moteur est associé à un capteur angulaire mesurant $\theta = G_t u_A$.

\subsubsection*{Relation entre $F_t$ et $G_t$}

On souhaite $F_t G_t = 1$, c'est-à-dire $F_t = 1/G_t$. Cette relation est vraie sous les hypothèses que la transmission est conservative (pas de pertes énergétiques), sans glissement, et en régime quasi-statique où la puissance est conservée : $C_m \dot{\theta} = f_A \dot{u_A}$.

\subsubsection*{Intérêt du jumeau numérique pour contrôler $u_B$}

Le système industriel ne dispose pas de mesure de $u_B$. Un jumeau numérique permet d'estimer $u_B$ à partir du modèle et des mesures disponibles (comme $u_A$ via $\theta$), afin de le contrôler en boucle fermée sans capteur direct.

\subsubsection*{Schéma blocs en boucle fermée avec mesure de $u_B$}

\begin{figure}[htbp]
    \centering
    \begin{tikzpicture}[auto, node distance=1.5cm]
        \node [input] (ref) {};
        \node [sum, right=0.5cm of ref] (sum) {};
        \node [block, right=1cm of sum] (F0) {$F_0$};
        \node [block, right=1cm of F0] (Ft) {$F_t$};
        \node [block, right=1cm of Ft] (F) {$F$};
        \node [coordinate, right=1cm of F] (split) {};

        % Sorties
        \node [block, above right=0.2cm and 1.5cm of split] (PiB) {$\Pi_B$};
        \node [block, below right=0.2cm and 1.5cm of split] (PiA) {$\Pi_A$};
        \node [output, right=1cm of PiB] (outB) {};
        \node [output, right=1cm of PiA] (outA) {};

        % Retour capteur
        \node [block, below=1cm of Ft] (C_inv) {$\hat{C}^{-1}$};
        \node [block, right=.75cm of C_inv] (C) {$C$};

        % Liaisons chaîne directe
        \draw [line] (ref) -- node[pos=1] {$+$} node[near start, above] {$e$} (sum);
        \draw [line] (sum) -- node {$I$} (F0); % e pour erreur ou signal commande
        \draw [line] (F0) -- node {$C_m$} (Ft);
        \draw [line] (Ft) -- node {$f_A$} (F);
        \draw [line] (F) -- node {$U$} (split);
        \draw [line] (split) |- (PiB);
        \draw [line] (split) |- (PiA);
        \draw [line] (PiB) -- node {$U_B$} (outB);
        \draw [line] (PiA) -- node {$U_A$} (outA);

        % Liaison de retour
        \draw [line] (PiB.east) -- ++(0.5,0) |- (C.east); % Prise sur UB
        \draw [line] (C) -- (C_inv);
        \draw [line] (C_inv) -| node[pos=0.95] {$-$} (sum);
    \end{tikzpicture}
    \caption{Schéma bloc du système en boucle fermée avec mesure de $U_B$}
    \label{fig:sys_boucle_fermee_mesure}
\end{figure}

On utilise un télémètre laser avec fonction de transfert $T$ pour mesurer $u_B$.

\subsubsection*{Schéma bloc du système asservi avec jumeau numérique}

\begin{figure}[!htbp]
    \centering
    \begin{tikzpicture}[auto, node distance=1.5cm]
        % Chaîne directe
        \node [input] (ref) {};
        \node [sum, right=0.5cm of ref] (sum) {};
        \node [block, right=0.8cm of sum] (K) {$K_{ctrl}$};
        \node [block, right=0.8cm of K] (F0) {$F_0$};
        \node [block, right=0.8cm of F0] (Ft) {$F_t$};
        \node [block, right=0.8cm of Ft] (F) {$F$};
        \node [coordinate, right=0.8cm of F] (split) {};
        
        \node [block, above right=0.2cm and 1cm of split] (PiB) {$\Pi_B$};
        \node [block, below right=0.2cm and 1cm of split] (PiA) {$\Pi_A$};
        \node [output, right=0.5cm of PiB] (outB) {};
        \node [output, right=0.5cm of PiA] (outA) {};

        % Chaîne Estimateur (Bas)
        % Positionnement relatif complexe pour aligner avec l'image
        \node [block, below=1.5cm of Ft] (F0_inv) {$F_0^{-1}$};
        \node [block, below=1.5cm of F0_inv] (Kalman) {$T_{kalman}$};
        \node [block, left=1cm of Kalman] (PiB_inv) {$\Pi_B^{-1}$}; % Notation de l'image
        
        \node [block, right=1cm of Kalman] (C_inv) {$\hat{C}^{-1}$};
        \node [block, right=1cm of C_inv] (C) {$C$};

        % --- TRACÉS ---
        
        % Chaîne principale
        \draw [line] (ref) -- node[pos=1] {$+$} node[near start, above] {$e$} (sum);
        \draw [line] (sum) -- node {$I$} (K);
        \draw [line] (sum) -- (K);
        \draw [line] (K) -- (F0);
        \draw [line] (F0) -- node[name=cm_node] {$C_m$} (Ft);
        \draw [line] (Ft) -- node {$f_A$} (F);
        \draw [line] (F) -- node {$U$} (split);
        \draw [line] (split) |- (PiB);
        \draw [line] (split) |- (PiA);
        \draw [line] (PiB) -- node[name=ub_node] {$U_B$} (outB);
        \draw [line] (PiA) -- node[name=ua_node] {$U_A$} (outA);

        % Entrées de l'estimateur (Kalman)
        % 1. Branche venant de Cm (commande)
        \draw [line] (cm_node.center) |- (F0_inv);
        \draw [line] (F0_inv) |- node[pos=0.7, right] {$f_A^{-1}$} (Kalman.130); % Entrée haut-gauche Kalman

        % 2. Branche venant de UA (mesure)
        \draw [line] (ua_node.south) |- (C.east);
        \draw [line] (C) -- (C_inv);
        \draw [line] (C_inv) -- node[above] {$\hat{U}_A$} (Kalman);

        % Sortie Estimateur et Feedback
        \draw [line] (Kalman) -- node[above] {$\hat{U}$} (PiB_inv);
        \draw [line] (PiB_inv) -| node[pos=0.05, above] {$\hat{U}_B$} node[pos=0.95] {$-$} (sum);
    \end{tikzpicture}
    \caption{Schéma bloc du système asservi avec jumeau numérique (filtre de Kalman)}
    \label{fig:sys_jumeau_numerique}
\end{figure}

On utilise un filtre de Kalman $T_{kalman}$ prenant la commande $C_m$ (après inversion par $F_0^{-1}$ pour obtenir $f_A^{-1}$) et la mesure filtrée $\hat{U}_A$ pour estimer $\hat{U}$, puis extraire $\hat{U}_B$ via $\Pi_B^{-1}$ pour le retour.

\end{document}